\begin{table}[!h]
\centering
\small
\caption{\label{tab:main}Main RDD Estimates: Effect of FPM Fiscal Windfalls on Homicide Rates}
\centering
\small
\begin{threeparttable}
\begin{tabular}[t]{lrrrr}
\toprule
Specification & Estimate & SE & $p$-value & Eff. N\\
\midrule
(1) Panel RDD (primary) & -0.286 & (2.473) & 0.903 & 62,865\\
(2) Log homicide rate & -0.0454 & (0.1488) & 0.887 & 62,865\\
(3) Panel + Year FE & -0.590 & (0.664) & 0.383 & 62,865\\
(4) Panel + State-Year FE & -0.573 & (0.551) & 0.308 & 62,865\\
(5) Youth homicides (15--29) & -0.139 & (1.335) & 0.931 & 62,865\\
\bottomrule
\end{tabular}
\begin{tablenotes}
\item \textit{Notes:} 
\item Dependent variable: homicide rate per 100,000 (rows 1, 3--4), log(rate + 1) (row 2), or youth (15--29) homicide rate (row 5). Row (1) reports the rdrobust estimate using annual population as the running variable with MSE-optimal bandwidth, triangular kernel, and robust bias-corrected inference. Rows (3)--(4) use local linear regression within the bandwidth with triangular kernel weights. Standard errors clustered at the state level. Annual panel, 2001--2021.
\end{tablenotes}
\end{threeparttable}
\end{table}
